% GENERATED FILE. Edit canonical lessons, then run npm run generate:books.
% canonical-source-hash: fnv1a64:72c166a64a56f761
% canonical-lessons: ML-C471-encin, ML-C471-cakram, ML-C471-takararu, ML-C471-hon, ML-C471-calunnuka

\chapter{At the Garage}
\label{ch:ml-at-the-garage}
\hlchaptermodality{471}

\begin{chapteropening}
\textbf{By the end of this chapter:} \emph{I can say an engine, a wheel, a fault, a horn and to get dented.}

Complete the last lesson of chapter 471: I can say an engine, a wheel, a fault, a horn and to get dented.
\end{chapteropening}

\section[\texorpdfstring{\ml{എഞ്ചിൻ} (eñcin)}{എഞ്ചിൻ (eñcin)}]{\ml{എഞ്ചിൻ} (eñcin) — an engine}

\label{lesson:ML-C471-encin}



\begin{quote}
Before the new one: say the Malayalam for a brake, then the Malayalam for to get a puncture.
\end{quote}

\subsection*{You'll want to know: \ml{എഞ്ചിൻ}}
\textbf{\ml{എഞ്ചിൻ}} — \emph{eñcin} — \textquotedblleft{}an engine\textquotedblright{}.

\begin{quote}
\ml{വണ്ടിയുടെ} \ml{എഞ്ചിൻ} \ml{വളരെ} \ml{പഴയതാണ്}.
\end{quote}

\begin{quote}
\emph{vaṇṭiyuṭe eñcin vaḷare paḻayatāṇŭ.}
\end{quote}

\begin{quote}
The vehicle's engine is very old.
\end{quote}

A train engine is also an \ml{എഞ്ചിൻ}.

\begin{cousinweb}
\textbf{\ml{എഞ്ചിൻ}} (\emph{eñcin}) was borrowed from English \emph{engine}, the machine that drives a vehicle.
\end{cousinweb}

\begin{grammarlens}[title={What you've built}]
One more word for getting around.
\end{grammarlens}

\subsection*{Your turn}
\begin{itemize}
\raggedright
  \item \emph{Say it:} \emph{eñcin}
  \item \emph{Say it:} \emph{eñcin}, once more
  \item \emph{Say it:} \emph{vaṇṭiyuṭe eñcin vaḷare paḻayatāṇŭ.} — The vehicle's engine is very old.
  \item \emph{Recall:} say the Malayalam for a brake, then the Malayalam for to get a puncture, then say \emph{eñcin} again
\end{itemize}

\begin{tcolorbox}[breakable,colback=teal!4,colframe=teal!35!black,title={Before you move on}]
What does \ml{എഞ്ചിൻ} mean? (\textquotedblleft{}An engine\textquotedblright{}.) Say it once more.
\end{tcolorbox}
\section[\texorpdfstring{\ml{ചക്രം} (cakraṁ)}{ചക്രം (cakraṁ)}]{\ml{ചക്രം} (cakraṁ) — a wheel}

\label{lesson:ML-C471-cakram}



\begin{quote}
Before the new one: say the Malayalam for to get a puncture, then the Malayalam for an engine.
\end{quote}

\subsection*{You'll want to know: \ml{ചക്രം}}
\textbf{\ml{ചക്രം}} — \emph{cakraṁ} — \textquotedblleft{}a wheel\textquotedblright{}.

\begin{quote}
\ml{ബസിന്റെ} \ml{ഒരു} \ml{ചക്രം} \ml{ചെളിയിൽ} \ml{താഴ്ന്നു}.
\end{quote}

\begin{quote}
\emph{basinṟe oru cakraṁ ceḷiyil tāḻnnu.}
\end{quote}

\begin{quote}
One of the bus's wheels sank into the mud.
\end{quote}

Plural: \ml{ചക്രങ്ങൾ}, \emph{cakraṅṅaḷ}.

A two-wheeler, a scooter or a motorbike, is an \ml{ഇരുചക്രവാഹനം}, a two-wheeled vehicle.

\begin{cousinweb}
\textbf{\ml{ചക്രം}} (\emph{cakraṁ}) is borrowed from Sanskrit \emph{cakra}, \textquotedblleft{}a wheel; a circle\textquotedblright{} (CDIAL 4538).
\end{cousinweb}

\begin{grammarlens}[title={What you've built}]
One more word for getting around.
\end{grammarlens}

\subsection*{Your turn}
\begin{itemize}
\raggedright
  \item \emph{Say it:} \emph{cakraṁ}
  \item \emph{Say it:} \emph{cakraṁ}, once more
  \item \emph{Say it:} \emph{basinṟe oru cakraṁ ceḷiyil tāḻnnu.} — One of the bus's wheels sank into the mud.
  \item \emph{Recall:} say the Malayalam for to get a puncture, then the Malayalam for an engine, then say \emph{cakraṁ} again
\end{itemize}

\begin{tcolorbox}[breakable,colback=teal!4,colframe=teal!35!black,title={Before you move on}]
What does \ml{ചക്രം} mean? (\textquotedblleft{}A wheel\textquotedblright{}.) Say it once more.
\end{tcolorbox}
\section[\texorpdfstring{\ml{തകരാറ്} (takarāṟŭ)}{തകരാറ് (takarāṟŭ)}]{\ml{തകരാറ്} (takarāṟŭ) — a fault; a breakdown}

\label{lesson:ML-C471-takararu}



\begin{quote}
Before the new one: say the Malayalam for an engine, then the Malayalam for a wheel.
\end{quote}

\subsection*{You'll want to know: \ml{തകരാറ്}}
\textbf{\ml{തകരാറ്}} — \emph{takarāṟŭ} — \textquotedblleft{}a fault; a breakdown\textquotedblright{}.

\begin{quote}
\ml{എഞ്ചിന്} \ml{എന്തോ} \ml{തകരാറുണ്ട്}.
\end{quote}

\begin{quote}
\emph{eñcinŭ entō takarāṟuṇṭŭ.}
\end{quote}

\begin{quote}
There is some fault in the engine.
\end{quote}

The thing with the fault takes -\ml{ന്} / -\ml{ക്ക്}: \ml{വണ്ടിക്ക്} \ml{തകരാറുണ്ട്}, the vehicle has a fault.

\begin{cousinweb}
\textbf{\ml{തകരാറ്}} (\emph{takarāṟŭ}) is usually traced to Arabic \emph{takrār}, \textquotedblleft{}repetition, dispute\textquotedblright{}, which came south through Hindustani; in Malayalam it names a fault in a machine.
\end{cousinweb}

\begin{grammarlens}[title={What you've built}]
One more word for getting around.
\end{grammarlens}

\subsection*{Your turn}
\begin{itemize}
\raggedright
  \item \emph{Say it:} \emph{takarāṟŭ}
  \item \emph{Say it:} \emph{takarāṟŭ}, once more
  \item \emph{Say it:} \emph{eñcinŭ entō takarāṟuṇṭŭ.} — There is some fault in the engine.
  \item \emph{Recall:} say the Malayalam for an engine, then the Malayalam for a wheel, then say \emph{takarāṟŭ} again
\end{itemize}

\begin{tcolorbox}[breakable,colback=teal!4,colframe=teal!35!black,title={Before you move on}]
What does \ml{തകരാറ്} mean? (\textquotedblleft{}A fault; a breakdown\textquotedblright{}.) Say it once more.
\end{tcolorbox}
\section[\texorpdfstring{\ml{ഹോൺ} (hōṇ)}{ഹോൺ (hōṇ)}]{\ml{ഹോൺ} (hōṇ) — a horn}

\label{lesson:ML-C471-hon}



\begin{quote}
Before the new one: say the Malayalam for a wheel, then the Malayalam for a fault.
\end{quote}

\subsection*{You'll want to know: \ml{ഹോൺ}}
\textbf{\ml{ഹോൺ}} — \emph{hōṇ} — \textquotedblleft{}a horn\textquotedblright{}.

\begin{quote}
\ml{പിന്നിലെ} \ml{കാർ} \ml{ഉറക്കെ} \ml{ഹോൺ} \ml{അടിച്ചു}.
\end{quote}

\begin{quote}
\emph{pinnile kār uṟakke hōṇ aṭiccu.}
\end{quote}

\begin{quote}
The car behind honked loudly.
\end{quote}

To sound a horn is \ml{ഹോൺ} \ml{അടിക്കുക}, to hit the horn.

\begin{cousinweb}
\textbf{\ml{ഹോൺ}} (\emph{hōṇ}) was borrowed from English \emph{horn}, the device a vehicle sounds as a warning.
\end{cousinweb}

\begin{grammarlens}[title={What you've built}]
One more word for getting around.
\end{grammarlens}

\subsection*{Your turn}
\begin{itemize}
\raggedright
  \item \emph{Say it:} \emph{hōṇ}
  \item \emph{Say it:} \emph{hōṇ}, once more
  \item \emph{Say it:} \emph{pinnile kār uṟakke hōṇ aṭiccu.} — The car behind honked loudly.
  \item \emph{Recall:} say the Malayalam for a wheel, then the Malayalam for a fault, then say \emph{hōṇ} again
\end{itemize}

\begin{tcolorbox}[breakable,colback=teal!4,colframe=teal!35!black,title={Before you move on}]
What does \ml{ഹോൺ} mean? (\textquotedblleft{}A horn\textquotedblright{}.) Say it once more.
\end{tcolorbox}
\section[\texorpdfstring{\ml{ചളുങ്ങുക} (caḷuṅṅuka)}{ചളുങ്ങുക (caḷuṅṅuka)}]{\ml{ചളുങ്ങുക} (caḷuṅṅuka) — to get dented}

\label{lesson:ML-C471-calunnuka}



\begin{quote}
Before the new one: say the Malayalam for a fault, then the Malayalam for a horn.
\end{quote}

\subsection*{You'll want to know: \ml{ചളുങ്ങുക}}
\textbf{\ml{ചളുങ്ങുക}} — \emph{caḷuṅṅuka} — \textquotedblleft{}to get dented\textquotedblright{}.

\begin{quote}
\ml{ചെറിയ} \ml{അപകടത്തിൽ} \ml{കാറിന്റെ} \ml{വാതിൽ} \ml{ചളുങ്ങി}.
\end{quote}

\begin{quote}
\emph{ceṟiya apakaṭattil kāṟinṟe vātil caḷuṅṅi.}
\end{quote}

\begin{quote}
In the small accident the car's door got dented.
\end{quote}

Past: \ml{ചളുങ്ങി}, \emph{caḷuṅṅi}.

Metal and tin things get dented: a pot, a car, a bucket.

\begin{cousinweb}
\textbf{\ml{ചളുങ്ങുക}} (\emph{caḷuṅṅuka}) is a Malayalam verb whose origin is not settled: DEDR 2407 lists a native \emph{caḷuṅṅuka}, \textquotedblleft{}to shake, to be frightened\textquotedblright{}, but whether the sense of being dented comes from that word is unclear.
\end{cousinweb}

\begin{grammarlens}[title={What you've built}]
One more word for getting around.
\end{grammarlens}

\subsection*{Your turn}
\begin{itemize}
\raggedright
  \item \emph{Say it:} \emph{caḷuṅṅuka}
  \item \emph{Say it:} \emph{caḷuṅṅuka}, once more
  \item \emph{Say it:} \emph{ceṟiya apakaṭattil kāṟinṟe vātil caḷuṅṅi.} — In the small accident the car's door got dented.
  \item \emph{Recall:} say the Malayalam for a fault, then the Malayalam for a horn, then say \emph{caḷuṅṅuka} again
\end{itemize}

\begin{tcolorbox}[breakable,colback=teal!4,colframe=teal!35!black,title={Before you move on}]
What does \ml{ചളുങ്ങുക} mean? (\textquotedblleft{}To get dented\textquotedblright{}.) Say it once more.
\end{tcolorbox}
